% Part: first-order-logic
% Chapter: sequent-calculus

\documentclass[../../../include/open-logic-chapter]{subfiles}

\begin{document}

\iftag{FOL}
      {\olchapter{fol}{seq}{De sequentencalculus}}
      {\olchapter{pl}{seq}{De sequentencalculus}}

\begin{editorial}
  Dit hoofdstuk behandelt Gentzens standaardsequentencalculus LK voor
  de klassieke eerste-ordelogica. Het kan nog met meer voorbeelden en
  opgaven worden aangevuld. Gebruik de tag ``prfLK'' om materiaal over
  de sequentencalculus als bewijssysteem op te nemen of uit te sluiten.
\end{editorial}

\olimport{rules-and-proofs}

\olimport{propositional-rules}

\iftag{FOL}{%
  \olimport{quantifier-rules}
}{}

\olimport{structural-rules}

\olimport{derivations}

\olimport{proving-things}

\iftag{FOL}{%
  \olimport{proving-things-quant}
}{}

\olimport{proof-theoretic-notions}

\olimport{provability-consistency}

\olimport{provability-propositional}

\iftag{FOL}{%
  \olimport{provability-quantifiers}
}{}

\olimport{soundness}

\iftag{FOL}{%
  \olimport{identity}
  \olimport{soundness-identity}
}{}

\OLEndChapterHook

\end{document}
